\chapter[Marco Metodológico]{Marco Metodológico}
\label{ch:metodologia}

\insertminitoc
\parindent0pt

El marco teórico desarrollado en los capítulos anteriores establece las bases conceptuales,
técnicas y normativas que sustentan el sistema propuesto. El presente capítulo describe
la metodología que se adoptará para su construcción y evaluación: el tipo y diseño de
investigación, los actores involucrados, las técnicas e instrumentos que se utilizarán para la
recolección y el procesamiento de los datos, el procedimiento que se seguirá a lo largo del
desarrollo y los métodos de validación que se emplearán.

\section{Tipo de investigación}

Por su propósito, la investigación se clasifica como aplicada \cite{sampieri_metodologia_2014}.
Su objetivo central no es generar conocimiento teórico nuevo, sino construir una
herramienta tecnológica funcional que resuelve un problema concreto: la brecha de
comunicación entre niños sordos, cuya lengua materna es el \acrshort{lesho}, y personas
oyentes que no conocen esa lengua. El punto de partida son conocimientos ya establecidos
en disciplinas como la visión por computadora, el aprendizaje profundo y el desarrollo
de aplicaciones móviles, que se combinan y adaptan para producir un sistema funcional
que facilita la comunicación entre ambas poblaciones en distintos contextos de la vida
cotidiana.

El enfoque adoptado es mixto \cite{sampieri_metodologia_2014}. La componente cuantitativa
abarca la construcción del conjunto de datos de landmarks, el entrenamiento y la
evaluación de los modelos de clasificación mediante métricas estándar como exactitud,
precisión, sensibilidad y puntuación F1, y la medición del tamaño de los modelos
exportados y los tiempos de inferencia en dispositivos reales. La componente cualitativa
corresponde a las sesiones de validación con usuarios de la Fundación CasAyuda,
donde se observa la interacción de los
participantes con la aplicación y se recopila retroalimentación sobre la comprensión de
la interfaz, la facilidad de uso y la adecuación de las señas reconocidas. Ambas
componentes se complementan: las métricas cuantitativas miden el desempeño técnico del
sistema, mientras que la observación cualitativa evalúa su pertinencia y usabilidad en
el contexto real para el que será diseñado.

El alcance de la investigación es exploratorio-descriptivo \cite{sampieri_metodologia_2014}.
Es exploratorio porque, al momento de realizar el estudio, no existen antecedentes
documentados de sistemas de reconocimiento automático de señas del \acrshort{lesho} entrenados
con técnicas de aprendizaje profundo, ni conjuntos de datos de landmarks de la lengua
disponibles públicamente. En ese sentido, la investigación abre un camino metodológico
en un dominio no explorado previamente en Honduras. Es descriptivo porque documenta de
forma sistemática el diseño del sistema, la arquitectura de los modelos, el proceso de
construcción del conjunto de datos y el comportamiento del sistema medido a través de
métricas reproducibles.

El diseño es no experimental \cite{sampieri_metodologia_2014}. No se conformará un grupo
de control ni se manipularán variables de forma aleatoria. El proceso seguirá un diseño de
investigación de desarrollo tecnológico: se definirán los requisitos del sistema, se
construirá el conjunto de datos, se entrenarán y validarán los modelos, se integrarán en una
aplicación móvil y se realizarán pruebas con usuarios reales. Este diseño es coherente con
la naturaleza del problema, donde el objetivo es demostrar la viabilidad técnica de una
solución antes de escalarla a un vocabulario más amplio o a condiciones de uso más
diversas.

\section{Diseño Metodológico}

La investigación se estructuró en cuatro fases secuenciales. La primera consistió en la
consulta con personas con conocimiento del \acrshort{lesho} para definir el vocabulario que el
sistema reconoce y muestra. La segunda comprendió la construcción del conjunto de datos
de landmarks mediante sesiones de grabación con cuatro participantes. La tercera abarcó
el entrenamiento de los modelos de clasificación y su integración en la aplicación móvil.
La cuarta fase corresponde a la validación del sistema con usuarios de la Fundación
CasAyuda. Cada fase produce entregables concretos que sirven como insumo para la
siguiente.

Una decisión metodológica transversal a todas las fases es el procesamiento
completamente local en el dispositivo. Esta decisión implica que ningún dato de cámara
se transmita a servidores externos, ni durante el uso de la aplicación ni durante la
captura del conjunto de datos. Los frames de video se procesarán en tiempo real con
MediaPipe Hands para extraer coordenadas de landmarks, y solo esas coordenadas se
almacenarán de forma permanente. Los videos crudos no se conservarán. Esta restricción,
motivada por consideraciones de privacidad, disponibilidad sin conexión a internet y
minimización de costos de infraestructura, condiciona las decisiones de diseño tanto del
script de captura de datos como de la arquitectura de la aplicación.

La elección de representar las señas mediante vectores de landmarks en lugar de imágenes
de video también responde a razones metodológicas. Al reducir cada mano a 21 puntos con
tres coordenadas, la entrada de los modelos es invariante ante variaciones de
iluminación, tono de piel y fondo, lo que hace que el conjunto de datos sea más
generalizable a distintas condiciones de uso, dispositivos y entornos de grabación.

\section{Población y Muestra}

\subsection{Población}

La población de esta investigación comprende tres grupos de personas involucrados en
situaciones de comunicación mediada por el \acrshort{lesho} en Honduras: niños y jóvenes
sordos que utilizan el \acrshort{lesho} como lengua materna, personas oyentes que interactúan
con ellos en contextos cotidianos, y personas con dominio del \acrshort{lesho} que participarán
como colaboradoras en la construcción del conjunto de datos y en la validación del
sistema. Estos tres grupos constituyen el universo de personas directamente implicadas
en el problema de comunicación que la aplicación busca atender.

\subsection{Muestra}

La muestra estará constituida por 30 personas seleccionadas mediante muestreo no
probabilístico intencional \cite{sampieri_metodologia_2014}. Este tipo de muestreo es
adecuado para investigaciones de carácter aplicado en las que no existe un marco
muestral que permita selección aleatoria y cuyo objetivo no es la generalización
estadística, sino la validación de un sistema funcional con participantes que cumplen
criterios específicos de inclusión.

El número de 30 participantes responde a tres criterios concurrentes. El primero es la
necesidad de representar los tres perfiles involucrados en el problema de comunicación:
personas sordas, personas con dominio del \acrshort{lesho} y personas oyentes. Sin presencia de
los tres grupos no sería posible evaluar ambas direcciones de comunicación del sistema
ni recopilar retroalimentación desde todas las perspectivas relevantes. El segundo
criterio es la suficiencia para el análisis descriptivo de los resultados del
cuestionario aplicado en la etapa de validación: muestras de esta magnitud permiten
calcular frecuencias, porcentajes y medidas de tendencia central de forma interpretable
sin requerir técnicas de inferencia estadística. El tercer criterio es la viabilidad
de acceso: la población de personas con dominio del \acrshort{lesho} en Honduras con afiliación
institucional verificable es reducida y está concentrada en organizaciones especializadas
como la Fundación CasAyuda, lo que limita de forma real el tamaño de muestra
alcanzable con garantías de idoneidad.

Los participantes se distribuirán en tres grupos según su función en la investigación:

\begin{itemize}
    \item \textbf{Personas sordas:} usuarios del módulo de reconocimiento de señas.
    Interactuarán con la aplicación ejecutando señas frente a la cámara y evaluarán la
    respuesta del sistema y la claridad de la interfaz.
    \item \textbf{Personas con dominio del \acrshort{lesho}:} participarán en la construcción del
    conjunto de datos ejecutando el vocabulario definido frente a la cámara, y también en
    las sesiones de validación evaluando la adecuación de las señas reconocidas y mostradas
    por el sistema.
    \item \textbf{Personas oyentes:} usuarios del módulo de texto a seña. Escribirán frases
    en español y evaluarán la representación visual generada por el sistema.
\end{itemize}

Los criterios de inclusión serán los siguientes: para el primer grupo, ser persona
sorda con uso activo del \acrshort{lesho}; para el segundo, tener dominio reconocido de la
lengua; para el tercero, ser persona oyente sin conocimiento previo del \acrshort{lesho}. No se
establecerá restricción por edad, género ni nivel de escolaridad, dado que la aplicación
se diseñará para entornos de uso variados.

\section{Técnicas e Instrumentos de Recolección de Datos}

La recolección de datos abarcará dos procesos diferenciados: la construcción del conjunto
de datos de landmarks para el entrenamiento de los modelos, y la validación del sistema
con usuarios reales. Cada proceso empleará técnicas e instrumentos distintos según sus
objetivos.

Para la construcción del conjunto de datos se utilizará la técnica de grabación controlada
de señas. Cada participante se ubicará frente a la cámara del dispositivo a una distancia
de entre 60 y 100 centímetros, con iluminación natural o interior estándar. El
instrumento principal será un script de captura desarrollado en Python que automatizará
todo el proceso: mostrará en pantalla la seña a ejecutar junto con una imagen de
referencia, presentará una cuenta regresiva de tres segundos, registrará la ejecución
durante dos a tres segundos y avanzará automáticamente a la siguiente seña sin
intervención manual. MediaPipe Hands se integrará directamente en
el script para extraer, en tiempo real, las coordenadas tridimensionales de los 21
landmarks de la mano detectada por cada frame. Dichas coordenadas se normalizarán
restando la posición de la muñeca a todos los demás puntos y se almacenarán en archivos
CSV etiquetados por clase. Los videos no se conservarán.

Para la validación del sistema con usuarios de la Fundación CasAyuda se
aplicará un cuestionario. El instrumento recogerá la percepción de los participantes
sobre la comprensión de la interfaz, la facilidad de interacción con la aplicación y la
adecuación de las señas reconocidas y mostradas por el sistema. Las respuestas
permitirán identificar aspectos de la experiencia de uso que no son evidentes a partir
de las métricas técnicas de los modelos.

\section{Procedimiento}

El desarrollo del sistema seguirá cinco fases secuenciales. Cada una producirá entregables
concretos que servirán como insumo para la siguiente etapa.

\subsection{Definición del vocabulario}

La primera fase consistirá en establecer con precisión las señas que el sistema
reconocerá y mostrará. En coordinación con personas con conocimiento del \acrshort{lesho}
vinculadas a la Fundación CasAyuda, se seleccionarán las 50 señas
dinámicas que conformarán el vocabulario del Modelo B. Los criterios de selección
priorizarán señas de uso cotidiano para niños, representativas de categorías como
emociones, familia, necesidades básicas, números y colores. En la misma etapa se definirán
dos señas de control, concebidas inicialmente para delimitar el inicio y el final de una
seña dinámica sin tocar la pantalla, que deberán ser visualmente distintas entre sí y de
cualquier letra del alfabeto dactilológico. Las 30 letras del alfabeto dactilológico,
incluidos los dígrafos CH, LL y RR, completarán el vocabulario del alfabeto, junto con una
clase de reposo. Al concluir esta fase, el vocabulario total del sistema quedará cerrado:
33 clases para el Modelo A y 50 señas dinámicas para el Modelo B.

\subsection{Construcción del conjunto de datos}

Con el vocabulario definido se realizaron las sesiones de grabación con cuatro
participantes, utilizando el script de captura descrito en la sección anterior.

El conjunto de datos del Modelo A abarca las 33 clases: las 30 letras del alfabeto
dactilológico, la clase REPOSO y las dos señas de control. Cada participante grabó 20
tomas de cada clase estática y 40 tomas de cada una de las cinco letras que se ejecutan con
movimiento (J, Ñ, Z, LL y RR), cuya trayectoria requiere más ejemplos. De cada toma se
extraen ventanas de 20 fotogramas, que constituyen la entrada del modelo; cada fotograma se
representa con 126 valores, correspondientes a las coordenadas $x$, $y$ y $z$ de los 21
landmarks de las dos manos, normalizadas respecto a la posición de la muñeca.

Para el conjunto de datos del Modelo B, cada participante grabó 25 tomas de cada una de
las 50 señas dinámicas, lo que conformó un total de 5\,000 secuencias con los
cuatro participantes. Cada secuencia comprende entre 30 y 60 fotogramas según la duración
natural de la seña, que se remuestrean a una longitud fija de 40. A diferencia del Modelo A,
cada fotograma del Modelo B combina la configuración de las manos (126 valores) con la
ubicación de estas respecto al cuerpo, obtenida con MediaPipe Pose (18 valores); con los
rasgos relativos que se derivan en el preprocesamiento, la entrada final al modelo es de 152
valores. Ambos conjuntos se almacenarán en formato CSV sin incluir video crudo. Como técnicas de aumento de datos se aplicarán variación de velocidad
mediante submuestreo y sobremuestreo temporal, y ruido gaussiano leve sobre las
coordenadas de los landmarks para simular la variación natural en la ejecución de cada
seña.

\subsection{Entrenamiento y exportación de los modelos}

Los conjuntos de datos se dividirán en subconjuntos de entrenamiento y evaluación para
el ajuste y validación de cada modelo.

El Modelo A se implementará como una red neuronal convolucional en una dimensión
(\acrshort{cnn} 1D) sobre el eje temporal de la ventana, con activación \acrshort{relu} y regularización
mediante dropout. La capa de salida utiliza la función softmax sobre 33 neuronas. El modelo se entrenará con el optimizador Adam y la
función de pérdida de entropía cruzada categórica, con detención temprana sobre el
conjunto de validación para prevenir el sobreajuste.

El Modelo B se implementará como una red \acrshort{lstm} en configuración many-to-one: la red
procesa la secuencia completa de frames y emite una única clasificación al final,
utilizando el estado oculto del último paso temporal como entrada a una capa densa de
salida softmax sobre 50 neuronas. Una vez entrenado cada modelo, se exportará al formato
\acrshort{tflite} mediante el conversor de TensorFlow Lite, con cuantización posterior al
entrenamiento para reducir el tamaño de los archivos y acelerar la inferencia en la
CPU del dispositivo.

\subsection{Desarrollo de la aplicación móvil}

Con los modelos exportados se procederá al desarrollo de la aplicación en Flutter. La
aplicación integra cuatro componentes principales: el módulo de cámara, el canal de
procesamiento de MediaPipe Hands, los intérpretes de los modelos \acrshort{tflite} y la
interfaz de usuario.

El módulo de reconocimiento de señas implementa el siguiente pipeline: cada frame
capturado por la cámara frontal se procesa con MediaPipe Hands para extraer los
landmarks, que se normalizan y se pasan al Modelo A. Las predicciones se filtran
mediante un mecanismo de persistencia temporal, que exige que la misma clase aparezca
durante al menos 5 ventanas consecutivas con una confianza igual o superior al 60\,\%,
y un período de enfriamiento de 1200 milisegundos tras cada detección confirmada. Cuando
la persona usuaria pulsa el botón de grabación, el sistema activa la captura de la
secuencia dinámica, que incorpora la ubicación en el cuerpo mediante MediaPipe Pose, y la
envía al Modelo B al pulsar de nuevo para detener. La extracción de landmarks con
MediaPipe se ejecuta en un hilo nativo en segundo plano, comunicado con la aplicación
mediante un canal de plataforma, para no bloquear el hilo principal de la interfaz.

El módulo de texto a seña implementa un buscador sobre un diccionario de clips de señas.
El usuario escribe en español y la aplicación reproduce la seña correspondiente a cada
palabra sobre un muñeco animado, dibujado en pantalla a partir de los landmarks del clip.
Para las palabras sin seña registrada, el sistema recurre al deletreo letra por letra con
los clips del alfabeto dactilológico.

\subsection{Validación con usuarios}

La validación del sistema se realizará en sesiones con usuarios de la Fundación CasAyuda.
Durante cada sesión, los participantes interactuarán con la
aplicación en ambas direcciones de comunicación. Al finalizar, se aplicará el
cuestionario descrito en la sección anterior. Los resultados, combinados con las
observaciones recogidas durante las sesiones, permitirán identificar aspectos de mejora
en la interfaz y confirmar la pertinencia del vocabulario incluido en el sistema.

\section{Herramientas Tecnológicas}

El desarrollo del proyecto requerirá dos conjuntos de herramientas diferenciados: uno
orientado a la captura del conjunto de datos y el entrenamiento de los modelos, y otro
orientado al desarrollo de la aplicación móvil.

\subsection{Herramientas para la captura de datos y el entrenamiento}

El script de captura y el pipeline de entrenamiento se desarrollarán en Python 3.10.
Las bibliotecas principales que se utilizarán son las siguientes:

\begin{itemize}
    \item \textbf{MediaPipe Hands}: para la extracción en tiempo real de los 21
    landmarks de la mano durante las sesiones de grabación.
    \item \textbf{TensorFlow y Keras}: para la definición, el entrenamiento y la
    exportación de los modelos convolucional 1D y \acrshort{lstm}, así como para la conversión al
    formato \acrshort{tflite} mediante el conversor incluido en el paquete.
    \item \textbf{NumPy}: para las operaciones matemáticas sobre los vectores de
    landmarks, incluyendo la normalización de coordenadas y la aplicación del aumento
    de datos.
    \item \textbf{Pandas}: para la lectura, organización y manipulación de los
    archivos CSV que componen el conjunto de datos.
    \item \textbf{Matplotlib}: para la generación de gráficos de entrenamiento y
    las matrices de confusión utilizadas en la evaluación de los modelos.
\end{itemize}

\subsection{Herramientas para el desarrollo de la aplicación}

La aplicación móvil se desarrollará con Flutter 3 en lenguaje Dart. Los paquetes
principales que se emplearán son los siguientes:

\begin{itemize}
    \item \textbf{camera}: para el acceso a la cámara frontal del dispositivo y la
    obtención del flujo de frames en tiempo real.
    \item \textbf{tflite\_flutter}: para cargar y ejecutar los modelos \acrshort{tflite}
    directamente en el dispositivo a través de la \acrshort{ffi}.
    \item \textbf{MediaPipe Hands y MediaPipe Pose}: integrados en la capa nativa de la
    aplicación para extraer los landmarks de las manos y, en las señas dinámicas, la
    ubicación en el cuerpo.
\end{itemize}

\subsection{Entorno de desarrollo y control de versiones}

Todo el desarrollo se realizará en Visual Studio Code con las extensiones oficiales de
Flutter y Python. El control de versiones se gestionará con Git, y el repositorio se
publicará en GitHub para documentar el progreso del proyecto y hacer accesible tanto el
código fuente como el conjunto de datos de landmarks.

\section{Métodos de Validación}

La validación del sistema se abordará desde dos dimensiones complementarias: la
evaluación técnica del desempeño de los modelos de clasificación y la evaluación de
la experiencia de uso con personas reales.

\subsection{Validación técnica de los modelos}

Ambos modelos se evaluarán sobre un subconjunto de datos reservado exclusivamente
para prueba, no utilizado durante el entrenamiento ni durante el ajuste de
hiperparámetros. El instrumento central de evaluación será la matriz de confusión, que
permite visualizar de forma detallada qué clases el modelo identifica correctamente y
con cuáles tiende a confundirse. A partir de la matriz de confusión se calcularán las
métricas estándar de clasificación: exactitud, precisión, sensibilidad y puntuación
F1. Dado que ambos modelos manejan múltiples clases con igual relevancia comunicativa,
se utilizará el macro-promedio para obtener una medida global que trate a todas las
clases por igual, independientemente de su frecuencia en el conjunto de datos. Las
métricas del Modelo A y el Modelo B se reportarán de forma independiente, ya que cada
uno resuelve un subproblema distinto del sistema.

\subsection{Validación con usuarios}

La validación con usuarios tendrá como objetivo evaluar aspectos del sistema que las
métricas técnicas no pueden capturar: la claridad de la interfaz, la facilidad de
interacción y la adecuación del vocabulario incluido. Para esto se realizarán sesiones
con usuarios de la Fundación CasAyuda, durante las cuales los
participantes interactuarán con la aplicación en ambas direcciones de comunicación.

Al finalizar cada sesión se aplicará el cuestionario descrito anteriormente. Las
respuestas se analizarán para identificar patrones en la percepción del sistema y
detectar aspectos que requieran ajuste antes de la versión final. Los resultados de
esta validación se presentarán junto con los de la evaluación técnica en el capítulo de
análisis de resultados.
